% GENERATED FILE. Edit canonical lessons, then run npm run generate:books.
% canonical-source-hash: fnv1a64:20ad386f880defeb
% canonical-lessons: RU-C199-molodyozh, RU-C199-grazhdanin, RU-C199-inostranets, RU-C199-turist, RU-C199-tekst, RU-R199-first-pass-words-from-chapters-191-194, RU-R199-second-pass-words-from-chapters-195-199

\chapter{Young people, Citizen, Foreigner, Tourist, Text}
\label{ch:ru-young-people-citizen-foreigner-tourist-text}
\hlchaptermodality{199}

\begin{chapteropening}
\textbf{By the end of this chapter:} \emph{I can say young people, a citizen, a foreigner, a tourist and a text.}

Complete the last lesson of chapter 199: I can say young people, a citizen, a foreigner, a tourist and a text.
\end{chapteropening}

\section[\texorpdfstring{\ru{молодёжь} (molodyózh)}{молодёжь (molodyózh)}]{\ru{молодёжь} (molodyózh) — young people}

\label{lesson:RU-C199-molodyozh}



\begin{quote}
Before the new one: say the Russian for a traffic jam; a cork, then the Russian for an accident.
\end{quote}

\subsection*{You'll want to know: \ru{молодёжь}}
\textbf{\ru{молодёжь}} (\emph{molodyózh}) — \textquotedblleft{}young people\textquotedblright{}.

\begin{grammarlens}[title={What you've built}]
One more word for this chapter.
\end{grammarlens}

\subsection*{Your turn}
\begin{itemize}
\raggedright
  \item \emph{Say it:} \emph{molodyózh}
  \item \emph{Say it:} \emph{molodyózh}, once more
  \item \emph{Say it:} say \emph{molodyózh}
  \item \emph{Recall:} say the Russian for a traffic jam; a cork, then the Russian for an accident, then say \emph{molodyózh} again
\end{itemize}

\begin{tcolorbox}[breakable,colback=teal!4,colframe=teal!35!black,title={Before you move on}]
What does \ru{молодёжь} mean? (\textquotedblleft{}Young people\textquotedblright{}.) Say it once more.
\end{tcolorbox}
\section[\texorpdfstring{\ru{гражданин} (grazhdanín)}{гражданин (grazhdanín)}]{\ru{гражданин} (grazhdanín) — a citizen}

\label{lesson:RU-C199-grazhdanin}



\begin{quote}
Before the new one: say the Russian for an accident, then the Russian for young people.
\end{quote}

\subsection*{You'll want to know: \ru{гражданин}}
\textbf{\ru{гражданин}} (\emph{grazhdanín}) — \textquotedblleft{}a citizen\textquotedblright{}.

\begin{grammarlens}[title={What you've built}]
One more word for this chapter.
\end{grammarlens}

\subsection*{Your turn}
\begin{itemize}
\raggedright
  \item \emph{Say it:} \emph{grazhdanín}
  \item \emph{Say it:} \emph{grazhdanín}, once more
  \item \emph{Say it:} say \emph{grazhdanín}
  \item \emph{Recall:} say the Russian for an accident, then the Russian for young people, then say \emph{grazhdanín} again
\end{itemize}

\begin{tcolorbox}[breakable,colback=teal!4,colframe=teal!35!black,title={Before you move on}]
What does \ru{гражданин} mean? (\textquotedblleft{}A citizen\textquotedblright{}.) Say it once more.
\end{tcolorbox}
\section[\texorpdfstring{\ru{иностранец} (inostránets)}{иностранец (inostránets)}]{\ru{иностранец} (inostránets) — a foreigner}

\label{lesson:RU-C199-inostranets}



\begin{quote}
Before the new one: say the Russian for young people, then the Russian for a citizen.
\end{quote}

\subsection*{You'll want to know: \ru{иностранец}}
\textbf{\ru{иностранец}} (\emph{inostránets}) — \textquotedblleft{}a foreigner\textquotedblright{}.

\begin{grammarlens}[title={What you've built}]
One more word for this chapter.
\end{grammarlens}

\subsection*{Your turn}
\begin{itemize}
\raggedright
  \item \emph{Say it:} \emph{inostránets}
  \item \emph{Say it:} \emph{inostránets}, once more
  \item \emph{Say it:} say \emph{inostránets}
  \item \emph{Recall:} say the Russian for young people, then the Russian for a citizen, then say \emph{inostránets} again
\end{itemize}

\begin{tcolorbox}[breakable,colback=teal!4,colframe=teal!35!black,title={Before you move on}]
What does \ru{иностранец} mean? (\textquotedblleft{}A foreigner\textquotedblright{}.) Say it once more.
\end{tcolorbox}
\section[\texorpdfstring{\ru{турист} (turíst)}{турист (turíst)}]{\ru{турист} (turíst) — a tourist}

\label{lesson:RU-C199-turist}



\begin{quote}
Before the new one: say the Russian for a citizen, then the Russian for a foreigner.
\end{quote}

\subsection*{You'll want to know: \ru{турист}}
\textbf{\ru{турист}} (\emph{turíst}) — \textquotedblleft{}a tourist\textquotedblright{}.

\begin{grammarlens}[title={What you've built}]
One more word for this chapter.
\end{grammarlens}

\subsection*{Your turn}
\begin{itemize}
\raggedright
  \item \emph{Say it:} \emph{turíst}
  \item \emph{Say it:} \emph{turíst}, once more
  \item \emph{Say it:} say \emph{turíst}
  \item \emph{Recall:} say the Russian for a citizen, then the Russian for a foreigner, then say \emph{turíst} again
\end{itemize}

\begin{tcolorbox}[breakable,colback=teal!4,colframe=teal!35!black,title={Before you move on}]
What does \ru{турист} mean? (\textquotedblleft{}A tourist\textquotedblright{}.) Say it once more.
\end{tcolorbox}
\section[\texorpdfstring{\ru{текст} (tekst)}{текст (tekst)}]{\ru{текст} (tekst) — a text}

\label{lesson:RU-C199-tekst}



\begin{quote}
Before the new one: say the Russian for a foreigner, then the Russian for a tourist.
\end{quote}

\subsection*{You'll want to know: \ru{текст}}
\textbf{\ru{текст}} (\emph{tekst}) — \textquotedblleft{}a text\textquotedblright{}.

\begin{grammarlens}[title={What you've built}]
One more word for this chapter.
\end{grammarlens}

\subsection*{Your turn}
\begin{itemize}
\raggedright
  \item \emph{Say it:} \emph{tekst}
  \item \emph{Say it:} \emph{tekst}, once more
  \item \emph{Say it:} say \emph{tekst}
  \item \emph{Recall:} say the Russian for a foreigner, then the Russian for a tourist, then say \emph{tekst} again
\end{itemize}

\begin{tcolorbox}[breakable,colback=teal!4,colframe=teal!35!black,title={Before you move on}]
What does \ru{текст} mean? (\textquotedblleft{}A text\textquotedblright{}.) Say it once more.
\end{tcolorbox}
\section[(dialogue)]{Words from chapters 191-194 — a first pass}

\label{lesson:RU-R199-first-pass-words-from-chapters-191-194}



\begin{quote}
Say the words for a tourist and a text.
\end{quote}

\subsection*{Your turn}
Say these aloud:

\begin{itemize}
\raggedright
  \item a radiator; a battery, a socket, a wire, a charger; morning exercises, a laptop
  \item a newspaper, a magazine, a dictionary, a ruler, a folder
  \item a brimmed hat, a belt, a pocket, a button, earrings
  \item dumplings, a salad, cabbage, beetroot, peas
\end{itemize}

\begin{tcolorbox}[breakable,colback=teal!4,colframe=teal!35!black,title={Before you move on}]
A radiator; a battery? (\textbf{\ru{батарея}}.) A socket? (\textbf{\ru{розетка}}.) A wire? (\textbf{\ru{провод}}.) A charger; morning exercises? (\textbf{\ru{зарядка}}.) A laptop? (\textbf{\ru{ноутбук}}.) A newspaper? (\textbf{\ru{газета}}.) A magazine? (\textbf{\ru{журнал}}.) A dictionary? (\textbf{\ru{словарь}}.) A ruler? (\textbf{\ru{линейка}}.) A folder? (\textbf{\ru{папка}}.) A brimmed hat? (\textbf{\ru{шляпа}}.) A belt? (\textbf{\ru{ремень}}.) A pocket? (\textbf{\ru{карман}}.) A button? (\textbf{\ru{пуговица}}.) Earrings? (\textbf{\ru{серьги}}.) Dumplings? (\textbf{\ru{пельмени}}.) A salad? (\textbf{\ru{салат}}.) Cabbage? (\textbf{\ru{капуста}}.) Beetroot? (\textbf{\ru{свёкла}}.) Peas? (\textbf{\ru{горох}}.)
\end{tcolorbox}
\section[(dialogue)]{Words from chapters 195-199 — a second pass}

\label{lesson:RU-R199-second-pass-words-from-chapters-195-199}



\begin{quote}
Say the words for nature and a branch.
\end{quote}

\subsection*{Your turn}
Say these aloud:

\begin{itemize}
\raggedright
  \item nature, a branch, a root, a rose, an oak
  \item a rat, a spider, a mosquito, a fly, a beetle
  \item a hobby, an exhibition, a play, a guitar, a piano
  \item a passenger, a carriage, a transfer, a traffic jam; a cork, an accident
  \item young people, a citizen, a foreigner, a tourist, a text
\end{itemize}

\begin{tcolorbox}[breakable,colback=teal!4,colframe=teal!35!black,title={Before you move on}]
Nature? (\textbf{\ru{природа}}.) A branch? (\textbf{\ru{ветка}}.) A root? (\textbf{\ru{корень}}.) A rose? (\textbf{\ru{роза}}.) An oak? (\textbf{\ru{дуб}}.) A rat? (\textbf{\ru{крыса}}.) A spider? (\textbf{\ru{паук}}.) A mosquito? (\textbf{\ru{комар}}.) A fly? (\textbf{\ru{муха}}.) A beetle? (\textbf{\ru{жук}}.) A hobby? (\textbf{\ru{хобби}}.) An exhibition? (\textbf{\ru{выставка}}.) A play? (\textbf{\ru{спектакль}}.) A guitar? (\textbf{\ru{гитара}}.) A piano? (\textbf{\ru{пианино}}.) A passenger? (\textbf{\ru{пассажир}}.) A carriage? (\textbf{\ru{вагон}}.) A transfer? (\textbf{\ru{пересадка}}.) A traffic jam; a cork? (\textbf{\ru{пробка}}.) An accident? (\textbf{\ru{авария}}.) Young people? (\textbf{\ru{молодёжь}}.) A citizen? (\textbf{\ru{гражданин}}.) A foreigner? (\textbf{\ru{иностранец}}.) A tourist? (\textbf{\ru{турист}}.) A text? (\textbf{\ru{текст}}.)
\end{tcolorbox}
